\section{Conclusão}\label{chap:conclusao-final}

O trabalho materializa uma plataforma computacional que reúne seleção geográfica, consulta a imagens, índices espectrais, séries temporais, contexto climático e exportação de resultados. A evolução para execução web preserva o SQLite local e adiciona isolamento de sessões e transferência GeoJSON, sem introduzir banco externo ou login de visitantes.

Os 677 testes aprovados, a cobertura integral dos módulos medidos com tracer C e a verificação mínima de acesso real ao Earth Engine constituem evidências de software, não validação agronômica. Em resposta à pergunta de pesquisa, a implementação demonstra uma forma de integrar os componentes e verificar seus contratos; a avaliação da plataforma como ferramenta de monitoramento agrícola ainda requer execução documentada de ponta a ponta e conferência independente dos cálculos.

Permanecem como próximos passos a validação da publicação no Streamlit Cloud, os cenários com culturas e geometrias confirmadas e as verificações de máscaras, resoluções e unidades climáticas. Até a obtenção dessas evidências, esta conclusão é parcial e não reivindica eficácia em campo, superioridade sobre outras ferramentas ou estudo de usabilidade concluído.

% PENDENTE: sintetizar o atendimento dos objetivos após a avaliação documentada.
% PENDENTE: apresentar contribuições efetivamente demonstradas, limitações e trabalhos futuros coerentes com os resultados.
